\documentclass[11pt,a4paper]{article}
\usepackage[margin=2.5cm]{geometry}
\usepackage[T1]{fontenc}
\usepackage{lmodern}
\usepackage{microtype}
\usepackage{enumitem}
\usepackage{needspace}
\usepackage{booktabs}
\usepackage{tabularx}
\usepackage[hidelinks]{hyperref}
\setlist{nosep,leftmargin=1.5em,topsep=3pt}
\setlength{\parskip}{6pt}
\setlength{\parindent}{0pt}
\newcolumntype{Y}{>{\raggedright\arraybackslash}X}

\newcommand{\point}[1]{\par\bigskip\needspace{6\baselineskip}\textbf{#1}\par}
\newcommand{\sub}[1]{\par\smallskip\needspace{4\baselineskip}\textit{#1}\par}

\begin{document}

\begin{center}
{\large\bfseries Response to the Reviewer}\\[6pt]
Previous title: \emph{BaitFL: Lightweight Trap-Based Free-Rider Detection and\\
Mitigation in Federated Learning}\\[3pt]
New title: \emph{SPARC-FL: Statistical Peer Auditing with Repeated\\
Confirmation for Free-Rider Detection in Federated Learning}
\end{center}

\bigskip
Dear Reviewer,

Thank you for the detailed review. It was very useful, and your first
comment in particular changed the direction of the paper. Below we first
explain the biggest change and give a short overview, then go through your
comments one by one. For each, we say what we changed, which experiments
are scheduled and why, and what we have not addressed. Section, table and
figure numbers refer to the revised manuscript.

% ---------------------------------------------------------------------
\point{1. We changed the title and the name of the method}

You pointed out that nothing in the paper showed the traps actually help.
We ran that test (details under M1), and they do not, at least for the
attacks and settings we studied. Sending the ordinary model in place of the
trap catches exactly the same free riders, and it does not throw away
about 11\% of honest training.

Since the paper could no longer honestly be called ``trap-based'', we
renamed the method SPARC-FL (Statistical Peer Auditing with Repeated
Confirmation) and changed the title. The contribution is now what the
evidence supports: comparing each client's update with peers, and removing
a client only when the same evidence comes back on two audit slots and in
a later cycle. The traps are kept only as an optional ``challenge
dispatch'' mode, reported together with their cost. The abstract,
introduction, contributions and Section IV were rewritten around this, and
Figures 2 and 3(a) were redrawn.

% ---------------------------------------------------------------------
\point{2. Overview}

\begin{tabularx}{\linewidth}{@{}l Y Y@{}}
\toprule
Comment & Addressed in this revision & Scheduled or not addressed \\
\midrule
M1 Traps & Ordinary versus challenge dispatch at 40\% (Section V-H) & Same comparison at 30\%. Zero-perturbation control. Norm-only, profile-only, send-back-only and no-memory variants \\
M2 Threat model & 5\%, 10\% and 20\% prevalence. Norm-matching diagnostic. Cohort minimum stated & Published decaying-noise and scaled-delta attacks. Full norm-matched FR4 and mixed attacks. S2-WEF not addressed \\
M3 Baseline & Failed STD-DAGMM setup disclosed. Norm threshold. Requirement-based comparison (Section II) & STD-DAGMM check in Lin et al.'s setting. No re-implementation of methods needing extra client protocol \\
M4 Cost & Absolute times, server-only share, withholding as its own cost & Synthetic scaling to 1,000 clients. Training gap disclosed, not explained \\
M5 Honest harm & Four honest-only runs, per-run uncertainty & $\alpha = 0.1$ and $1.0$. Additional seeds \\
M6 References & Rebuilt and extended & -- \\
Minor 1--6 & All addressed & -- \\
\bottomrule
\end{tabularx}

% ---------------------------------------------------------------------
\point{M1. The traps are never shown to matter}

\sub{What we did}
We ran the ablation you asked for. The ordinary-dispatch version keeps
everything the same (audit schedule, response slots, norm and profile
statistics, send-back check and removal rule) but sends the normal global
model and aggregates the responses. We matched it against the original
runs on both datasets, both partitions and all four attacks with 40\%
free riders (Section V-H, Table VI, Figure 6).

\sub{What we found}
Both versions removed all 640 free riders and no honest client. They
finished in the same round in 15 of the 16 settings, and the ordinary
version was faster in the last one (round 50 against 60). Accuracy
differed by at most 0.3 percentage points. The ordinary version withheld
at most 0.7\% of honest client-rounds, compared with 11\% to 18.5\% with
traps.

We also tried the simple norm threshold you suggested, a one-round robust
screen using the median and MAD. It flagged no honest client but caught all
free riders in only 5 of the 16 settings, because with so many fabricated
updates the median and MAD are pulled towards them. So the detection comes
from the per-slot peer comparison and the repeated-confirmation rule, not
from the traps and not from a single norm cut-off. You were also right that
the additive FR4 response cancels the trap, and Section V-E now says so.

\sub{Scheduled}
\begin{itemize}
\item The same comparison at 30\% free riders, so it covers both
  prevalences of the main study.
\item A control that keeps the frozen dispatched model but sets the
  perturbation to zero. It separates the effect of freezing the model
  within a cycle from the effect of the noise itself.
\item The four variants you listed: norm-only, profile-only,
  send-back-only and no retained memory. In the norm-only and profile-only
  variants the other statistic is switched off before nomination, not just
  ignored afterwards, so they test each statistic on its own.
\end{itemize}

% ---------------------------------------------------------------------
\point{M2. The threat model is too weak}

\sub{What we did}
\begin{itemize}
\item \textbf{Lower prevalence.} We ran 5\%, 10\% and 20\% free riders for
  all four attacks on both datasets and partitions (Section V-G, Table V).
\item \textbf{Norm matching.} We replayed all 72 initial FR2 to FR4
  decisions from the main study with each free rider's norm replaced by the
  honest median, which gives the attacker extra information (Section V-F,
  Table IV).
\item \textbf{Collusion.} Section VII-D now says that colluders can compare
  hashes of the models they receive, and that copying another client's
  update cancels a trap just as FR4 does. We make no collusion claim.
\end{itemize}

\sub{What we found}
Your concern was correct. At 5\%, FR2 to FR4 were not caught in any
setting, and only FR1 was removed through the exact send-back check. At
10\% and 20\% all free riders were removed, with one honest removal across
all of these runs. The misses at 5\% come from the cohort-support rule,
and the abstract now states that statistical removal needs a nominated
group of at least ten clients when 100 are active. One small note: the
rule limits the size of the nominated group rather than the number of free
riders, so a few free riders could in principle be caught inside a larger
mixed group. That did not happen in our runs.

Under norm matching, the norm test caught no attacker. The profile test
still caught 61.9\% of FR2 cases, but only 10.7\% of FR3 and 17.9\% of FR4,
with ten honest false certificates. We report this openly as a weakness.

\sub{Scheduled}
\begin{itemize}
\item The decaying-noise attack of Fraboni et al.\ (the disguised free
  rider) and the scaled-delta attack of Zhu et al. These are the published
  attacks closest to your concern and do not rely on small updates.
\item A full-model norm-matched FR4, trained end to end rather than replayed
  on stored features, where the attacker is given the honest median norm
  of the previous round. This is deliberately stronger than a real attacker.
\item Mixed strategies, with FR1 to FR4 present in the same run.
\end{itemize}

\sub{Not addressed}
The S2-WEF adaptive attack, which depends on the WEF client protocol
discussed under M3.

% ---------------------------------------------------------------------
\point{M3. There is no working detection baseline}

\sub{What we did}
We agree that a baseline which flags no free rider is not a useful
comparison. The paper now says plainly that our five-epoch STD-DAGMM setup
did not work as a detector, keeps our chosen training budget separate from
the original method, and makes no claim of better detection. The timing
comparison is presented only as a measurement of that setup (Section V-K).
We added WEF-defense, S2-WEF, FLET and FLoomChecker to the related work,
and the simple norm threshold described under M1.

\sub{Why we did not re-implement the other detectors}
SPARC-FL is meant for a specific setting: the server sees only the returned
models, does not change what clients train on, does not ask clients for
extra reports, and keeps its own computation small. Most of the suggested
detectors need something outside this setting, so a matched benchmark would
mean changing the threat model to fit each of them. Section II now sets out
these requirements:
\begin{itemize}
\item \textbf{STD-DAGMM} fits an autoencoder to full update vectors every
  round. With the encoder widths of our reproduction, its first layer alone
  has 128 weights per model coordinate, about
  26.5 million for our MNIST model and 79.5 million for CIFAR-10, so the
  detector is far larger than the model it audits. Fitting took 15.76 s and
  42.09 s per round in our runs, against 0.215 s and 0.791 s for our
  checking.
\item \textbf{FRIDA} needs honest clients to train on server-supplied
  canaries, which changes their local objective, or infers properties of
  their private data.
\item \textbf{FRAD} relies on client-reported resource and contribution
  information, and \textbf{PASS} on clients evaluating each other's models.
\item \textbf{WEF-defense and S2-WEF} need clients to compute and upload
  weight-evolution matrices, which the server cannot check against the
  returned model. S2-WEF also simulates attack patterns on the server.
\item \textbf{FLET and FLoomChecker} add test tasks or integrity evidence on
  the client side.
\end{itemize}
For the cosine threshold you mentioned, Section V-H gives an argument rather
than a measurement. FR4 adds the average of recent global changes, so its
direction follows the global trend, and FR1's zero update has no direction.
A cosine cut-off strict enough to catch FR4 would also flag non-IID honest
clients.

\sub{Scheduled}
A check of STD-DAGMM in the setting of Lin et al.\ (MNIST, 100 clients, the
published two-layer model and attacker), to see whether the method works
when its own conditions are met.

\sub{Not addressed}
A matched detection benchmark against the methods above. We accept that this
leaves the comparison partly argued rather than measured, and the paper says
so in Section VII-E.

% ---------------------------------------------------------------------
\point{M4. The ``lightweight'' claim uses a misleading denominator}

\sub{What we did}
\begin{itemize}
\item The abstract now gives absolute checking times, 215 ms per round on
  MNIST and 791 ms on CIFAR-10, and compares them with server-side work only
  (about 88\% and 89\% of checking plus aggregation), not with the whole
  simulation.
\item ``Lightweight'' is defined once, as measured server computation at
  100 clients (Section V-I).
\item Withheld honest training is reported as a separate cost (Section
  V-D). As M1 shows, ordinary dispatch removes almost all of it.
\item Section V-K explains that client diagnostic evaluation is only
  simulator bookkeeping, is not needed in deployment and is left out of
  checking time. The 8.95 s CIFAR-10 training gap is reported as
  unexplained, since our records cannot tell us the cause.
\end{itemize}

\sub{Scheduled}
\begin{itemize}
\item A scaling test of the leave-one-out step at 100, 500 and 1,000
  clients, timed separately for feature construction and decisions. It uses
  synthetic client features, so it measures the server step alone.
\item A reference run that excludes the free riders from the start, with no
  auditing. It separates the accuracy gained by removing free riders from
  the accuracy lost by withholding honest updates.
\end{itemize}

\sub{Not addressed}
A controlled rerun that would explain the CIFAR-10 training-time gap.

% ---------------------------------------------------------------------
\point{M5. Honest-client harm is under-tested}

\sub{What we did}
We added four runs with no free riders, on both datasets and partitions.
None removed an honest client (Section V-G). Four runs is a small number,
and we say so: the exact 95\% interval for the chance of any honest removal
in a run is $[0, 0.60]$. We do not treat the 6,240 honest client-cases in
the main study as independent, because clients in one federation share
training and reference statistics. ``Unresolved'' is kept only as a
defined label, not as an outcome, and all status terms are defined once in
Table X.

\sub{Scheduled}
\begin{itemize}
\item Stronger and weaker non-IID data ($\alpha = 0.1$ and $1.0$), each with
  an honest-only run and runs with FR3 and FR4. False removals are expected
  to be most likely at $\alpha = 0.1$.
\item Repetitions of the main settings with new random seeds, to measure
  variation between runs properly, and the remaining matched no-defense
  runs so every accuracy comparison has a partner.
\end{itemize}

% ---------------------------------------------------------------------
\point{M6. Reference and positioning problems}

All reference links were rebuilt from the bibliography and checked in the
compiled PDF. WEF-defense, S2-WEF (marked as a preprint), FLET and
FLoomChecker were added to Section II. We cite WEF-defense in the version
whose details we could verify. These works are discussed as related
approaches and were not run as baselines.

% ---------------------------------------------------------------------
\point{Minor comments}
\begin{enumerate}
\item The abstract now states the minimum group size, the withheld
  training under traps and that the attackers are non-adaptive.
\item Client diagnostic evaluation is explained in Section V-K and is not
  needed in deployment.
\item ``Lightweight'' means measured server computation at $N = 100$.
\item The equations and Table I were rebuilt directly in LaTeX and checked
  in the PDF.
\item The captions of the score plots (Figure 3(b) and Figures 11 and 12)
  now say they show averaged coordinates, not the per-slot tests.
\item Evidence outcomes and participation states are defined once in
  Table X and used consistently.
\end{enumerate}

\par\bigskip\needspace{9\baselineskip}
Thank you again for the time you put into this review. It led to the main
result of this revision, even if it was not the result we were hoping for,
and the paper is more accurate because of it.

\par\bigskip
Best regards,\\
Aashnan Rahman

\end{document}
